\documentclass[10pt,a4paper]{amsart}
\usepackage{amsmath,amssymb,amsthm,graphicx}
% Retained minimal preamble for the Stacks Project chapter more-algebra.tex (More on Algebra).
% The chapter file pinned-original-source.tex carries only \input{preamble} on its first line, so no
% preamble text lives inside the chapter itself.  The Stacks Project uses one single shared preamble for
% every chapter; the snapshot stages it identically as sources/stacks-chow/preamble.tex and
% sources/stacks-etale-cohomology/preamble.tex.  The block below is copied mechanically, line for line,
% from sources/stacks-chow/preamble.tex lines 14-18 (the Xy-pic setup) and 154-217 (the native theorem
% declarations together with the mathematical macros).  Nothing was added, removed or redefined; the
% wrapper class and the amsmath/amssymb/amsthm/graphicx base packages are supplied by the assembler.
% Chapter: more-algebra.tex -- More on Algebra.
% For commutative diagrams we use Xy-pic
\usepackage[all]{xy}

% We use 2cell for 2-commutative diagrams.
\xyoption{2cell}

% Theorem environments.
%
\theoremstyle{plain}
\newtheorem{theorem}[subsection]{Theorem}
\newtheorem{proposition}[subsection]{Proposition}
\newtheorem{lemma}[subsection]{Lemma}

\theoremstyle{definition}
\newtheorem{definition}[subsection]{Definition}
\newtheorem{example}[subsection]{Example}
\newtheorem{exercise}[subsection]{Exercise}
\newtheorem{situation}[subsection]{Situation}

\theoremstyle{remark}
\newtheorem{remark}[subsection]{Remark}
\newtheorem{remarks}[subsection]{Remarks}

\numberwithin{equation}{subsection}

% Macros
%
\def\lim{\mathop{\mathrm{lim}}\nolimits}
\def\colim{\mathop{\mathrm{colim}}\nolimits}
\def\Spec{\mathop{\mathrm{Spec}}}
\def\Hom{\mathop{\mathrm{Hom}}\nolimits}
\def\Ext{\mathop{\mathrm{Ext}}\nolimits}
\def\SheafHom{\mathop{\mathcal{H}\!\mathit{om}}\nolimits}
\def\SheafExt{\mathop{\mathcal{E}\!\mathit{xt}}\nolimits}
\def\Sch{\mathit{Sch}}
\def\Mor{\mathop{\mathrm{Mor}}\nolimits}
\def\Ob{\mathop{\mathrm{Ob}}\nolimits}
\def\Sh{\mathop{\mathit{Sh}}\nolimits}
\def\NL{\mathop{N\!L}\nolimits}
\def\CH{\mathop{\mathrm{CH}}\nolimits}
\def\proetale{{pro\text{-}\acute{e}tale}}
\def\etale{{\acute{e}tale}}
\def\QCoh{\mathit{QCoh}}
\def\Ker{\mathop{\mathrm{Ker}}}
\def\Im{\mathop{\mathrm{Im}}}
\def\Coker{\mathop{\mathrm{Coker}}}
\def\Coim{\mathop{\mathrm{Coim}}}

% Boxtimes
%
\DeclareMathSymbol{\boxtimes}{\mathbin}{AMSa}{"02}

%
% Macros for moduli stacks/spaces
%
\def\QCohstack{\mathcal{QC}\!\mathit{oh}}
\def\Cohstack{\mathcal{C}\!\mathit{oh}}
\def\Spacesstack{\mathcal{S}\!\mathit{paces}}
\def\Quotfunctor{\mathrm{Quot}}
\def\Hilbfunctor{\mathrm{Hilb}}
\def\Curvesstack{\mathcal{C}\!\mathit{urves}}
\def\Polarizedstack{\mathcal{P}\!\mathit{olarized}}
\def\Complexesstack{\mathcal{C}\!\mathit{omplexes}}
% \Pic is the operator that assigns to X its picard group, usage \Pic(X)
% \Picardstack_{X/B} denotes the Picard stack of X over B
% \Picardfunctor_{X/B} denotes the Picard functor of X over B
\def\Pic{\mathop{\mathrm{Pic}}\nolimits}
\def\Picardstack{\mathcal{P}\!\mathit{ic}}
\def\Picardfunctor{\mathrm{Pic}}
\def\Deformationcategory{\mathcal{D}\!\mathit{ef}}

% --- appended 2026-10-12 (Chief): the project own definitions of the reference/slogan/history
% display environments, copied mechanically from the shared Stacks preamble lines 8-12
% (\usepackage{verbatim} plus the three \newenvironment lines, which the chapter preamble keeps
% verbatim from the project preamble). Nothing added, removed or redefined: these environments
% carry no mathematical content and the project itself typesets them as comments. Without these
% three lines a counted statement whose own text uses them fails with Environment slogan undefined.
\usepackage{verbatim}
\newenvironment{reference}{\comment}{\endcomment}
\newenvironment{slogan}{\comment}{\endcomment}
\newenvironment{history}{\comment}{\endcomment}
\usepackage{cleveref}
\crefname{originallemma}{original lemma}{original lemmas}
\Crefname{originallemma}{Original lemma}{Original lemmas}
\makeatletter
\newlabel{more-algebra-lemma-formally-smooth-JZ}{{0DYH}{source}{}{sourcepointer-1}{}}
\newlabel{more-algebra-lemma-formally-smooth-JZ@cref}{{[originallemma][0][]0DYH}{[1][0][]source}{}{}{}}
\makeatother
\crefname{originallemma}{original lemma}{original lemmas}
\Crefname{originallemma}{Original lemma}{Original lemmas}
\makeatletter
\newlabel{algebra-lemma-NL-surjection}{{07BP}{source}{}{sourcepointer-2}{}}
\newlabel{algebra-lemma-NL-surjection@cref}{{[originallemma][0][]07BP}{[1][0][]source}{}{}{}}
\makeatother
\crefname{originallemma}{original lemma}{original lemmas}
\Crefname{originallemma}{Original lemma}{Original lemmas}
\makeatletter
\newlabel{algebra-lemma-exact-sequence-NL}{{00S2}{source}{}{sourcepointer-3}{}}
\newlabel{algebra-lemma-exact-sequence-NL@cref}{{[originallemma][0][]00S2}{[1][0][]source}{}{}{}}
\makeatother
\crefname{originallemma}{original lemma}{original lemmas}
\Crefname{originallemma}{Original lemma}{Original lemmas}
\makeatletter
\newlabel{lemma-conormal-sequence-H1-regular}{{0666}{source}{}{sourcepointer-4}{}}
\newlabel{lemma-conormal-sequence-H1-regular@cref}{{[originallemma][0][]0666}{[1][0][]source}{}{}{}}
\makeatother
\crefname{originallemma}{original lemma}{original lemmas}
\Crefname{originallemma}{Original lemma}{Original lemmas}
\makeatletter
\newlabel{lemma-noetherian-finite-all-equivalent}{{09CC}{source}{}{sourcepointer-5}{}}
\newlabel{lemma-noetherian-finite-all-equivalent@cref}{{[originallemma][0][]09CC}{[1][0][]source}{}{}{}}
\makeatother
\crefname{originallemma}{original lemma}{original lemmas}
\Crefname{originallemma}{Original lemma}{Original lemmas}
\makeatletter
\newlabel{lemma-formally-smooth-flat}{{07NP}{source}{}{sourcepointer-6}{}}
\newlabel{lemma-formally-smooth-flat@cref}{{[originallemma][0][]07NP}{[1][0][]source}{}{}{}}
\makeatother
\crefname{originallemma}{original lemma}{original lemmas}
\Crefname{originallemma}{Original lemma}{Original lemmas}
\makeatletter
\newlabel{theorem-regular-fs}{{07EL}{source}{}{sourcepointer-7}{}}
\newlabel{theorem-regular-fs@cref}{{[originallemma][0][]07EL}{[1][0][]source}{}{}{}}
\makeatother
\crefname{originallemma}{original lemma}{original lemmas}
\Crefname{originallemma}{Original lemma}{Original lemmas}
\makeatletter
\newlabel{lemma-base-change-fs}{{07EG}{source}{}{sourcepointer-8}{}}
\newlabel{lemma-base-change-fs@cref}{{[originallemma][0][]07EG}{[1][0][]source}{}{}{}}
\makeatother
\title{The Jacobi-Zariski sequence of a formally smooth local homomorphism at the residue field}
\author{The Stacks Project Authors}
\date{stacks-project:more-algebra.tex:local-snapshot-20261012; GNU FDL 1.2 or later}
\begin{document}
\maketitle
\setcounter{subsection}{1}
\noindent\textit{Editorial note.} All delivered mathematical text is the author's own, copied from the pinned Stacks Project chapter file more-algebra.tex; no statement or proof step was rewritten, reordered or removed.\par\medskip
\noindent\textit{Editorial note.} The chapter file pinned-original-source.tex carries only \textbackslash{}input\{preamble\} on its first line, so no preamble text lives inside the chapter. The retained file sources/stacks-more-algebra/retained-preamble.tex is a mechanical, line-for-line copy of the project's single shared preamble as staged in the snapshot (Xy-pic setup plus theorem declarations and mathematical macros). Nothing was added, removed or redefined, and the excerpt serves only as preamble, never as a span coordinate.\par\medskip
\noindent\textit{Editorial note.} Both counted spans start at the \textbackslash{}begin line of their environment, so the statement and the author's proof are delivered whole; surrounding prose of the section is deliberately outside the counted spans and is carried as standard prerequisites instead.\par\medskip
\noindent\textit{Editorial note.} The citation fields (source\_id, source\_version, authors, source\_url, license.name, source\_title) are identical to those of every other card this lane has produced from this chapter, so the batch selector treats them as one source.\par\medskip
\noindent\textit{Allowed standard prerequisites (editorial).}\begin{itemize}
\item Identification of the cotangent space with H\_1 of the naive cotangent complex, and exactness of the Jacobi-Zariski sequence. Algebra, Lemma \textbackslash{}ref\{algebra-lemma-NL-surjection\} and Algebra, Lemma \textbackslash{}ref\{algebra-lemma-exact-sequence-NL\}, cited at proof lines 10029-10031. Reduce the claim to injectivity of H\_1(NL\_\{K/A\}) into m\_B/m\_B\textasciicircum{}2.
\item The conormal sequence for a regular quotient and the criterion for a finite module to be cut by a regular sequence. Lemma \textbackslash{}ref\{lemma-conormal-sequence-H1-regular\} and Lemma \textbackslash{}ref\{lemma-noetherian-finite-all-equivalent\}, cited at proof lines 10048-10050. Produce the short exact sequence 0 to m\_A/m\_A\textasciicircum{}2 tensor K to m\_B/m\_B\textasciicircum{}2 to m\_\{B-bar\}/m\_\{B-bar\}\textasciicircum{}2 to 0.
\item Formal smoothness of a local map implies flatness, and formal smoothness passes to the special fibre. Lemma \textbackslash{}ref\{lemma-formally-smooth-flat\} and Lemma \textbackslash{}ref\{lemma-base-change-fs\}, cited at proof lines 10052 and 10072. Give flatness of A to B and formal smoothness of the fibre extension, which is what makes the left vertical arrow of the diagram injective.
\item Injectivity of the map on differentials induced by a formally smooth fibre extension. Theorem \textbackslash{}ref\{theorem-regular-fs\}, cited at proof line 10071. Supplies the injectivity of the left vertical arrow, after which the snake lemma completes the proof.
\end{itemize}\medskip
\ifdefined\hypertarget\hypertarget{sourcepointer-1}{}\fi
\noindent\textit{Source reference.} Original lemma 0DYH: focused-seeds/AGP/sources/stacks-snapshot-20261011/tags.txt, source lines 17467--17467. The lemma's own canonical tag, quoted so that readers can resolve the label inside the delivered text; it points at the counted statement itself and carries no proof content.\par\medskip
\ifdefined\hypertarget\hypertarget{sourcepointer-2}{}\fi
\noindent\textit{Source reference.} Original lemma 07BP: focused-seeds/AGP/sources/stacks-snapshot-20261011/tags.txt, source lines 9349--9349. Navigation pointer to a statement the counted proof refers to but does not contain; nothing load-bearing is delegated to it.\par\medskip
\ifdefined\hypertarget\hypertarget{sourcepointer-3}{}\fi
\noindent\textit{Source reference.} Original lemma 00S2: focused-seeds/AGP/sources/stacks-snapshot-20261011/tags.txt, source lines 1024--1024. Navigation pointer to a statement the counted proof refers to but does not contain; nothing load-bearing is delegated to it.\par\medskip
\ifdefined\hypertarget\hypertarget{sourcepointer-4}{}\fi
\noindent\textit{Source reference.} Original lemma 0666: focused-seeds/AGP/sources/stacks-snapshot-20261011/tags.txt, source lines 7925--7925. Navigation pointer to a statement the counted proof refers to but does not contain; nothing load-bearing is delegated to it.\par\medskip
\ifdefined\hypertarget\hypertarget{sourcepointer-5}{}\fi
\noindent\textit{Source reference.} Original lemma 09CC: focused-seeds/AGP/sources/stacks-snapshot-20261011/tags.txt, source lines 11825--11825. Navigation pointer to a statement the counted proof refers to but does not contain; nothing load-bearing is delegated to it.\par\medskip
\ifdefined\hypertarget\hypertarget{sourcepointer-6}{}\fi
\noindent\textit{Source reference.} Original lemma 07NP: focused-seeds/AGP/sources/stacks-snapshot-20261011/tags.txt, source lines 9769--9769. Navigation pointer to a statement the counted proof refers to but does not contain; nothing load-bearing is delegated to it.\par\medskip
\ifdefined\hypertarget\hypertarget{sourcepointer-7}{}\fi
\noindent\textit{Source reference.} Original lemma 07EL: focused-seeds/AGP/sources/stacks-snapshot-20261011/tags.txt, source lines 9451--9451. Navigation pointer to a statement the counted proof refers to but does not contain; nothing load-bearing is delegated to it.\par\medskip
\ifdefined\hypertarget\hypertarget{sourcepointer-8}{}\fi
\noindent\textit{Source reference.} Original lemma 07EG: focused-seeds/AGP/sources/stacks-snapshot-20261011/tags.txt, source lines 9446--9446. Navigation pointer to a statement the counted proof refers to but does not contain; nothing load-bearing is delegated to it.\par\medskip
\begin{lemma}
\label{lemma-formally-smooth-JZ}
Let $A \to B$ be a local homomorphism of Noetherian local rings.
Assume $A \to B$ is formally smooth in the $\mathfrak m_B$-adic
topology. Let $K$ be the residue field of $B$. Then
the Jacobi-Zariski sequence for $A \to B \to K$ gives an exact sequence
$$
0 \to H_1(\NL_{K/A}) \to \mathfrak m_B/\mathfrak m_B^2
\to \Omega_{B/A} \otimes_B K \to \Omega_{K/A} \to 0
$$
\end{lemma}

\begin{proof}
Observe that $\mathfrak m_B/\mathfrak m_B^2 = H_1(\NL_{K/B})$
by Algebra, Lemma \ref{algebra-lemma-NL-surjection}.
By Algebra, Lemma \ref{algebra-lemma-exact-sequence-NL}
it remains to show injectivity of
$H_1(\NL_{K/A}) \to \mathfrak m_B/\mathfrak m_B^2$.
With $k$ the residue field of $A$, the Jacobi-Zariski sequence
for $A \to k \to K$ gives $\Omega_{K/A} = \Omega_{K/k}$ and
an exact sequence
$$
\mathfrak m_A/\mathfrak m_A^2 \otimes_k K \to
H_1(\NL_{K/A}) \to
H_1(\NL_{K/k}) \to 0
$$
Set $\overline{B} = B \otimes_A k$. Since $\overline{B}$ is
regular the ideal $\mathfrak m_{\overline{B}}$ is generated
by a regular sequence. Applying
Lemmas \ref{lemma-conormal-sequence-H1-regular} and
\ref{lemma-noetherian-finite-all-equivalent}
to $\mathfrak m_A B \subset \mathfrak m_B$
we find $\mathfrak m_A B / (\mathfrak m_AB \cap \mathfrak m_B^2) =
\mathfrak m_A B / \mathfrak m_A \mathfrak m_B$ which is
equal to $\mathfrak m_A/\mathfrak m_A^2 \otimes_k K$
as $A \to B$ is flat by Lemma \ref{lemma-formally-smooth-flat}.
Thus we obtain a short exact sequence
$$
0 \to
\mathfrak m_A/\mathfrak m_A^2 \otimes_k K \to
\mathfrak m_B/\mathfrak m_B^2 \to
\mathfrak m_{\overline{B}}/\mathfrak m_{\overline{B}}^2 \to 0
$$
Functoriality of the Jacobi-Zariski sequences shows that
we obtain a commutative diagram
$$
\xymatrix{
&
\mathfrak m_A/\mathfrak m_A^2 \otimes_k K \ar[d] \ar[r] &
H_1(\NL_{K/A}) \ar[d] \ar[r] &
H_1(\NL_{K/k}) \ar[d] \ar[r] & 0 \\
0 \ar[r] &
\mathfrak m_A/\mathfrak m_A^2 \otimes_k K \ar[r] &
\mathfrak m_B/\mathfrak m_B^2 \ar[r] &
\mathfrak m_{\overline{B}}/\mathfrak m_{\overline{B}}^2 \ar[r] & 0
}
$$
The left vertical arrow is injective by Theorem \ref{theorem-regular-fs}
as $k \to \overline{B}$ is formally smooth in
the $\mathfrak m_{\overline{B}}$-adic topology by
Lemma \ref{lemma-base-change-fs}.
This finishes the proof by the snake lemma.
\end{proof}

\end{document}
